\documentclass[12pt]{article}
\usepackage{hyperref}

\usepackage{graphicx}

\usepackage{mathtools}
\DeclarePairedDelimiter\ceil{\lceil}{\rceil}
\DeclarePairedDelimiter\floor{\lfloor}{\rfloor}

\usepackage{cancel}

\usepackage{amsmath}
\usepackage{amsfonts}
\usepackage{amssymb}

\usepackage{listings}
\lstset{language=C++,
                basicstyle=\ttfamily,
                keywordstyle=\color{blue}\ttfamily,
                stringstyle=\color{red}\ttfamily,
                commentstyle=\color{green}\ttfamily,
                morecomment=[l][\color{magenta}]{\#}
}

\usepackage{breqn}
\usepackage[]{algorithm2e}
\usepackage{physics}

\newcommand\fnurl[2]{%
  \href{#2}{#1}\footnote{\url{#2}}%
}

\newcommand{\Four}{\mathcal{F}}
\renewcommand{\(}{\left(}
\renewcommand{\)}{\right)}
%% \renewcommand{\{}{\left{}
%% \renewcommand{\}}{\right}}


\title{Hermitian Symmetry and Real/Complex Transforms}
\author{Malcolm Roberts}



\def\no#1{\nobreak{#1}} % stop equation breaking in dmath

\begin{document}
\maketitle

\section{Introduction}

Hermitian symmetry: how important is it?

\subsection{Properties of real/complex DFTs}

The DFT is defined as
\begin{dmath}
  F_k = \sum_{n=0}^{N-1} f_n \omega_N^{n k}
\end{dmath}
where $\omega_N^{n k} = e^{-i 2 \pi/N}$.  If the input
$\{f_n\}_{n=0}^{N-1}$ is real-valued, then
\begin{dmath}
  F_{N-k}
  = \sum_{n=0}^{N-1} f_n \omega_N^{n (N-k)}
  = \sum_{n=0}^{N-1} f_n \,\cancelto{1}{\omega_N^{n N}}\,  \omega_N^{-nk}
  = \sum_{n=0}^{N-1} f_n  \omega_N^{-nk}
  = \(\sum_{n=0}^{N-1} f_n  \omega_N^{nk}\)^*
  = F_{k}^*
\end{dmath}
where $a^*$ denotes the complex conjugate of $a$.  This relationship
wihere $k\rightarrow N-k$ is equivalent to complex conjugation is
called Hermitian symmetry. This holds for the inverse as well: $f$ is
real valued if and only if $F$ is Hermitian-symmetric.

This gives rise to the real/complex Fourier transform, which is more
efficient than doing the complex version.  By using the symmetry of
$F$, we only need to store $\floor{N/2}+1$ complex values.  Moreover,
by taking advantage of the cyclic nature of the roots of unity, we can
see that
\begin{dmath}
  F_0^* = F_{N-0} = F_0
\end{dmath}
which implies that $F_0$, being its own complex-conjugate, has zero
imaginary part.  If $2 \mid N$, we also have
\begin{dmath}
  F_{N/2}^* = F_{N-N/2} = F_{N/2}
\end{dmath}
so $F_{N/2}$ (referred to as the Nyquist frequency) is also purely
real-valued.

The simplest way to compute a complex-to-real transform is to expand
the input from lengtt $\floor{N/2} +1$ to $N$ by making use of the
symmetry, performing an inverse complex FFT, and then taking the real
part.  We will use this algorithm as our reference computation.

For 2D transforms of size $N_x \times N_y$, the symmetry is
\begin{dmath}
  F_{i,j}^* = F_{N_x-i, N_y-j}
\end{dmath}
We could reduce symmetrize in either the $x$ or $y$ direction; the
convention is that the last direction is used for symmetry.  Thus, we
end up with a complex array of dimension $N_x \times \floor{N_y/2}+1$.
Depending on the parity of $N_x$ and $N_y$ we have either one, two, or
four values which are purely real.  In addition,
\begin{dmath}
  F_{i,0}^*
  = F_{N_x-i, N_y-0}
  = F_{N_x-i, 0}
\end{dmath}
so the the $y$-axis is itself Hermitian.  This is also true of the
$y$-Nyquist axis if $2\mid{}N_x$.

For 3D data, we have 1, 2, 4, or 8 purely real values at the corners,
one or two $x$-axes, one or two $y$-axes, and one or two $xy$-planes
which have Hermitian symmetry.

When computing the inverse transform, the symmetrized dimension is the
last to be transformed.  The other dimension transforms will produce
an output which is purely real whenever a sub-dimensional array is
Hermitian symmetric.  This implies, for example, that the batched 1D
complex-to-real will always have zero imaginary part in the DC mode
(and the Nyquist mode if even-length).  Thus, it is not necessary to
read in this data.  However, having the conjugate-symmetry in place
is necessary.

This is shown in \texttt{hermitian.py}.  When we compute a
multi-dimensional 3D complex-to-real transform on random complex data
of size $N_x\times{}N_y\times{}N_z$ the output of the complex-embedded
transform does not equal the output of the complex-embedded transform
if Hermitian symmetry is imposed before the transform.  However, if we
only impose the conjugate-symmetry, but not the zero imaginary part at
the corners, we still get the same result.

\end{document}
